\chapter{Conditional Probability and Independence}

\textcolor{red}{\textit{In all exercises the probability space is fixed, and $A$, $B$, $A_n$, etc... are events.}}

\begin{problem}
  Show that if $A \cap B = \varnothing$, then $A$ and $B$ \textit{cannot} be independent unless $P(A) = 0$ or $P(B) = 0$.
\end{problem}

\begin{solution}
  Suppose $P(A) = 0$ or $P(B) = 0$. Then, $P(A \cap B) = P(\varnothing) = 0 = P(A)P(B)$ and so $A$ and $B$ are independent. If neither $P(A)$ nor $P(B)$ is zero, then $P(A)P(B) \neq 0$, but $P(A \cap B) = 0$, and so $P(A \cap B) \neq P(A)P(B)$, i.e., $A$ and $B$ are not independent.
\end{solution}

%====================================================================================================

\begin{problem}
  Let $P(C) > 0$. Show that $P(A \cup B | C) = P(A | C) + P (B | C) - P(A \cap B | C)$.
\end{problem}

\begin{solution}
  We have 
  \begin{align*}
  P(A \cup B | C) &= \frac{P((A \cup B) \cap C)}{P(C)} \\
  &= \frac{P((A \cap C) \cup (B \cap C))}{P(C)} \\
  &= \frac{P(A \cap C) + P(B \cap C) - P((A \cap C) \cap (B \cap C))}{P(C)} \\
  &= \frac{P(A \cap C)}{P(C)} + \frac{P(B \cap C)}{P(C)} - \frac{P(A \cap B \cap C)}{P(C)} \\
  &= P(A|C) + P(B|C) - P(A \cap B | C),
  \end{align*}
  as desired.
\end{solution}

%====================================================================================================

\begin{problem}
  Suppose $P(C) > 0$ and $A_1, ..., A_n$ are all pairwise disjoint. Show that 
  \[ P\left( \bigcup_{i = 1}^n A_i | C \right) = \sum_{i = 1}^n P(A_i | C). \]
\end{problem}

\begin{solution}
  Since $A_1, ..., A_n$ are all pairwise disjoint, we have
  \begin{align*}
  P \left( \bigsqcup_{i = 1}^n A_i | C \right) &= \frac{P \left( \bigsqcup_{i = 1}^n (A_i \cap C) \right)}{P(C)} \\
  &= \frac{\sum_{i = 1}^{n} P(A_i \cap C)}{P(C)} \\
  &= \sum_{i = 1}^n \frac{P(A_i \cap C)}{P(C)} \\
  &= \sum_{i = 1}^n P(A_i | C),
  \end{align*}
  as desired.
\end{solution}

%====================================================================================================

\begin{problem}
  Let $P(B) > 0$. Show that $P(A \cap B) = P(A | B)P(B)$.
\end{problem}

\begin{solution}
  Since $P(B) > 0$, the conditional probability of $A$ given $B$, by definition, is given by $P(A | B) = P(A \cap B) / P(B)$, and so $P(A \cap B) = P(A | B) P(B)$.
\end{solution}

%====================================================================================================

\begin{problem}
  Let $0 < P(B) < 1$ and let $A$ be any event. Show
  \[ P(A) = P(A | B) P(B) + P(A | B^c) P(B^c). \]
\end{problem}

\begin{solution}
  Since $B$ and $B^c$ form a partition of the state space $\Omega$, by the Partition Equation, $P(A) = P(A | B) P(B) + P(A | B^c) P(B^c)$.
\end{solution}

%====================================================================================================

\begin{problem}
  Donated blood is screened for AIDS. Suppose the test has 99% accuracy, and that one in ten thousand people in your age group are HIV positive. The test has 5% false positive rating as well. Suppose the test screens you as positive. What is the probability you have AIDS? Is it 99%? (\textit{Hint}: 99% refers to $P$(test positive | you have AIDS). You want to find $P$(you have AIDS | test is positive).)
\end{problem}

\begin{solution}
  Let $A$ denote the event that you have AIDS, and $T$ denote the event that the test is positive. Then, $P(A) = \frac{1}{10000}$, $P(T | A) = \frac{99}{100}$, and $P(T | A^c) = \frac{5}{100}$. We want to find $P(A | T)$. By Bayes' Theorem, $P(A | T) = \frac{P(T | A) P(A)}{P(T)}$, and by Problem 5,
  \[ P(A | T) = \frac{P(T | A) P(A)}{P(T | A) P(A) + P(T | A^c) P(A^c)} = \frac{0.99 \cdot 0.0001}{0.99 \cdot 0.0001 + 0.05 \cdot 0.9999} \approx 0.00198. \]
\end{solution}
